% 第 5 章 Working With Designs in Memory
\bichapter{处理内存中的设计}{Working With Designs in Memory}
\label{chap:memory}

DC Explorer 将设计从设计文件读入内存。任意时刻内存中都可以同时存在多个设计。
设计读入后，可用多种方式修改，例如对子设计进行 group/ungroup，或更改子设计引用。

要了解如何打开与修改设计，见：
\begin{itemize}
  \item 设计术语
  \item 打开设计
  \item 设置当前设计
  \item 链接设计
  \item 编辑设计
  \item 更改设计层次
  \item 使用设计对象属性
  \item 创建设计
  \item 复制设计
  \item 重命名设计
  \item 保存设计
  \item 移除设计
  \item 设计数据库与网表命名一致性
\end{itemize}

% ============================================================
\section{设计术语}
\subsection*{Design Terminology}

不同公司对设计及其组件使用不同术语。
本节描述 Synopsys 综合工具使用的术语，以及设计实例与引用之间的关系。

本节包括：
\begin{itemize}
  \item 设计
  \item 设计对象
  \item 设计、实例与引用之间的关系
\end{itemize}

\subsection{设计}
\subsubsection*{Designs}

设计是执行逻辑功能的电路描述。
设计以多种设计格式描述，如 VHDL 或 Verilog HDL。
逻辑级设计表示为布尔方程组；门级设计（如网表）表示为互连单元。
设计可彼此独立综合，也可用作更大设计中的子设计。

设计可以是平坦的或层次化的：
\begin{itemize}
  \item \textbf{平坦设计（Flat designs）}\\
        平坦设计不包含子设计，只有一层结构，仅包含库单元。
  \item \textbf{层次化设计（Hierarchical designs）}\\
        层次化设计包含一个或多个设计作为子设计。
        每个子设计可进一步包含子设计，从而形成多层设计层次。
        包含子设计的设计称为父设计（parent designs）。
\end{itemize}

\subsection{设计对象}
\subsubsection*{Design Objects}

图~\ref{fig:dce-design-objs} 给出名为 TOP 的设计中的设计对象。
Synopsys 命令、属性与约束可指向特定设计对象。

\figplaceholder{Figure 19: Design Objects}{设计对象}{fig:dce-design-objs}

\subsubsection{设计}
\subsubsection*{Designs}

设计由实例、网络、端口与引脚组成；可包含子设计与库单元。
在图~\ref{fig:dce-design-objs} 中，设计为 TOP、ENCODER 与 REGFILE。
活动设计（正在处理的设计）称为当前设计（current design）。
多数命令特定于当前设计，即在当前设计的上下文中操作。

\subsubsection{引用}
\subsubsection*{References}

引用（reference）是可用作构建更大电路之元素的库组件或设计。
引用的结构可以是简单逻辑门，也可以是更复杂的设计（如 RAM 核或 CPU）。
一个设计可包含同一引用的多次出现；每次出现是一个实例。

reference 使您可以在某一设计内优化每个 cell（如 NAND 门），而不影响其他设计中的 cell。
某一设计中的引用独立于不同设计中的相同引用。
在图~\ref{fig:dce-design-objs} 中，引用为 INV、ENCODER 与 REGFILE。

\subsubsection{实例或单元}
\subsubsection*{Instances or Cells}

instance 是已加载到内存的 reference（库组件或设计）在电路中的一次实例化；每个 instance 都有唯一名称。
一个设计可包含多个实例。
多个实例可指向同一引用，但每个实例有唯一名称以与其他实例区分。
实例也称为单元（cell）。

设计中另一设计的唯一实例称为层次实例（hierarchical instance）。
设计中库单元的唯一实例称为叶单元（leaf cell）。
某些命令在当前设计的层次实例上下文中工作。
当前实例（current instance）为这些实例相关命令定义活动实例。
在图~\ref{fig:dce-design-objs} 中，实例为 U1、U2、U3 与 U4。

\subsubsection{端口}
\subsubsection*{Port}

端口是设计的输入与输出。端口方向指定为 input、output 或 inout。

\subsubsection{引脚}
\subsubsection*{Pins}

引脚是设计中单元（如门与触发器）的输入与输出。
子设计的端口在父设计中表现为引脚。

\subsubsection{网络}
\subsubsection*{Nets}

网络是将端口连接到引脚、以及将引脚彼此连接的连线。

\subsection{设计、实例与引用之间的关系}
\subsubsection*{Relationships Between Designs, Instances, and References}

图~\ref{fig:dce-inst-ref} 给出设计、实例与引用之间的关系。

\figplaceholder{Figure 20: Instances and References}{实例与引用}{fig:dce-inst-ref}

EXREF 设计包含两个引用：NAND2 与 MULTIPLIER。
NAND2 被实例化三次，MULTIPLIER 被实例化一次。
赋给 NAND2 三个实例的名称分别为 U1、U2 与 U3。
EXREF 设计中 NAND2 与 MULTIPLIER 的引用独立于不同设计中的相同引用。

\begin{seeAlsoBox}
\begin{itemize}
  \item 解析引用的过程
\end{itemize}
\end{seeAlsoBox}

% ============================================================
\section{打开设计}
\subsection*{Opening Designs}

打开设计之前，请确保设计采用下表所列格式之一。

\begin{table}[htbp]
\centering
\caption{支持的输入格式}
\label{tab:dce-input-fmt}
\begin{tabular}{@{}ll@{}}
\toprule
格式 & 说明 \\
\midrule
\texttt{.ddc} & Synopsys 内部数据库格式（推荐） \\
SystemVerilog & IEEE 标准 SystemVerilog（见 HDL Compiler 文档） \\
Verilog & IEEE 标准 Verilog（见 HDL Compiler 文档） \\
VHDL & IEEE 标准 VHDL（见 HDL Compiler 文档） \\
\bottomrule
\end{tabular}
\end{table}

要打开设计，使用下列两种方法之一：
\begin{itemize}
  \item \cmd{analyze} 与 \cmd{elaborate} 命令\\
        例如：
\begin{lstlisting}
de_shell> analyze -format vhdl my_design
de_shell> elaborate my_design
\end{lstlisting}
        若 \cmd{analyze} 报告错误，请在 HDL 源文件中修复后再次运行 \cmd{analyze}。
        修改已分析的设计后，必须重新分析。
  \item \cmd{read\_file} 命令\\
        例如：
\begin{lstlisting}
de_shell> read_file -rtl vhdl my_file
\end{lstlisting}
\end{itemize}

要了解两种方法的差异以及如何打开不同格式的设计，见：
\begin{itemize}
  \item 两种读入方法的差异
  \item 打开 HDL 与 \texttt{.ddc} 文件
\end{itemize}

\subsection{\texttt{analyze} 与 \texttt{elaborate} 命令}
\subsubsection*{The analyze and elaborate Commands}

\cmd{analyze} 命令执行下列任务：
\begin{itemize}
  \item 读入 HDL 源文件
  \item 在不为设计构建逻辑的情况下检查错误
  \item 以与 HDL 无关的中间格式创建 HDL 库对象
  \item 将中间文件存储在您定义的位置
\end{itemize}

\cmd{elaborate} 命令执行下列任务：
\begin{itemize}
  \item 使用通用技术（GTECH）库元素将中间设计转换为与工艺无关的逻辑设计
  \item 允许更改源码中定义的参数值
  \item 允许选择 VHDL architecture
  \item 用 DesignWare 组件替换代码中的 HDL 算术运算符
  \item 自动执行 \cmd{link} 命令以解析设计引用
\end{itemize}

更多信息亦见：
\begin{itemize}
  \item \textit{The HDL Compiler for Verilog User Guide}
  \item \textit{The HDL Compiler for VHDL User Guide}
\end{itemize}

\subsection{\texttt{read\_file} 命令}
\subsubsection*{The read\_file Command}

\cmd{read\_file} 命令：
\begin{itemize}
  \item 可读入多种不同格式
  \item 在单一步骤中执行与 \cmd{analyze} 和 \cmd{elaborate} 相同的操作
  \item 为 VHDL 创建 \texttt{.mr} 与 \texttt{.st} 中间文件
  \item 不自动执行 \cmd{link} 命令
  \item 不为 Verilog 创建任何中间文件
\end{itemize}

要创建中间文件，将 \varname{hdlin\_auto\_save\_templates} 变量设为 \texttt{true}。

对内存中的设计，DC Explorer 使用 \texttt{path\_name/design.ddc} 命名约定。
\texttt{path\_name} 是原文件读入所在目录，\texttt{design} 是设计名。
若随后从同名文件读入设计，DC Explorer 会覆盖原设计。
为防止覆盖，对 \cmd{read\_file} 使用 \opt{-single\_file} 选项。

若未指定设计格式，\cmd{read\_file} 根据文件扩展名推断格式。
若未使用已知扩展名，DC Explorer 使用 \texttt{.ddc} 格式。
用于自动推断的受支持扩展名不区分大小写。
\cmd{read\_file} 可对除 \texttt{.ddc} 外的所有格式读入压缩文件（\texttt{.ddc} 写出时已在内部压缩）。
要让 DC Explorer 自动推断压缩文件的格式，使用下列文件命名结构：\texttt{filename.format.gz}。

支持的扩展名包括：
\begin{itemize}
  \item \texttt{ddc}：\texttt{.ddc} 扩展名
  \item \texttt{db}：\texttt{.db}、\texttt{.sldb}、\texttt{.sdb}、\texttt{.db.gz}、\texttt{.sldb.gz} 与 \texttt{.sdb.gz} 扩展名
  \item Verilog：\texttt{.v}、\texttt{.verilog}、\texttt{.v.gz} 与 \texttt{.verilog.gz} 扩展名
  \item SystemVerilog：\texttt{.sv}、\texttt{.sverilog}、\texttt{.sv.gz} 与 \texttt{.sverilog.gz} 扩展名
  \item VHDL：\texttt{.vhd}、\texttt{.vhdl}、\texttt{.vhd.gz} 与 \texttt{.vhdl.gz} 扩展名
\end{itemize}

\begin{seeAlsoBox}
\begin{itemize}
  \item 链接设计
\end{itemize}
\end{seeAlsoBox}

\subsection{两种读入方法的差异}
\subsubsection*{Differences Between the Two Read Methods}

表~\ref{tab:dce-read-diff} 汇总使用 \cmd{analyze}/\cmd{elaborate} 与使用 \cmd{read\_file} 的差异。

\begin{table}[htbp]
\centering
\caption{\cmd{analyze}/\cmd{elaborate} 与 \cmd{read\_file} 的对比}
\label{tab:dce-read-diff}
\small
\begin{tabular}{@{}p{2.2cm}p{5.5cm}p{5.5cm}@{}}
\toprule
对比项 & \cmd{analyze} 与 \cmd{elaborate} & \cmd{read\_file} \\
\midrule
输入格式 & 支持 VHDL 与 Verilog 格式。 & 支持所有格式。 \\
何时使用 & 综合 VHDL 或 Verilog 文件。 & 综合网表、预编译设计等。 \\
参数 & 允许在 \cmd{elaborate} 命令行设置参数值。对于 parameterized design，可用 \cmd{analyze}/\cmd{elaborate} 以非默认值构建新设计。 & 不能传递参数；必须在 HDL 中使用 directive。 \\
Architecture & 允许指定要 elaborate 的 architecture。 & 不能指定要 elaborate 的 architecture。 \\
链接过程 & \cmd{elaborate} 自动执行 \cmd{link} 以解析引用。 & 必须使用 \cmd{link} 命令解析引用。 \\
\bottomrule
\end{tabular}
\end{table}

更多信息亦见：
\begin{itemize}
  \item \textit{The HDL Compiler for Verilog User Guide}
  \item \textit{The HDL Compiler for VHDL User Guide}
\end{itemize}

\subsection{打开 HDL 与 \texttt{.ddc} 文件}
\subsubsection*{Opening HDL and .ddc Files}

可打开 HDL 与 \texttt{.ddc} 格式的设计：
\begin{itemize}
  \item 打开 HDL 文件
  \item 打开 \texttt{.ddc} 文件
\end{itemize}

\subsubsection{打开 HDL 文件}
\subsubsection*{Opening HDL Files}

要打开 HDL 文件，使用下列方法之一：
\begin{itemize}
  \item \cmd{analyze} 与 \cmd{elaborate} 命令\\
        使用此方法时，按自底向上顺序分析顶层设计与所有子设计，然后 elaborate 顶层设计以及需要赋值或覆盖参数的任何子设计。\\
        例如：
\begin{lstlisting}
de_shell> analyze -format vhdl -library -work RISCTYPES.vhd
de_shell> analyze -format vhdl -library -work {ALU.vhd STACK_TOP.vhd \
  STACK_MEM.vhd ...}
de_shell> elaborate RISC_CORE \
  -architecture STRUCT -library WORK -update
\end{lstlisting}
  \item \cmd{read\_file} 命令\\
        例如：
\begin{lstlisting}
de_shell> read_file -format verilog RISC_CORE.v
\end{lstlisting}
  \item \cmd{read\_verilog} 或 \cmd{read\_vhdl} 命令\\
        例如：
\begin{lstlisting}
de_shell> read_verilog RISC_CORE.v
\end{lstlisting}
        也可使用 \cmd{read\_file} \opt{-format} \texttt{VHDL} 或 \cmd{read\_file} \opt{-format} \texttt{verilog} 命令。
\end{itemize}

\subsubsection{打开 \texttt{.ddc} 文件}
\subsubsection*{Opening .ddc Files}

要打开 \texttt{.ddc} 文件，使用 \cmd{read\_ddc} 命令或 \cmd{read\_file} \opt{-format} \texttt{ddc} 命令。例如：
\begin{lstlisting}
de_shell> read_ddc RISC_CORE.ddc
\end{lstlisting}

\texttt{.ddc} 格式向后兼容但不向前兼容；即，可读取用较旧版本 DC Explorer 创建的 \texttt{.ddc} 文件，但不能读取比当前所用版本更新的版本创建的文件。

% ============================================================
\section{设置当前设计}
\subsection*{Setting the Current Design}

要将正在处理的设计设为当前设计，使用下列命令之一：
\begin{itemize}
  \item \cmd{read\_file} 命令\\
        \cmd{read\_file} 成功完成后，将当前设计设为刚读入的设计。例如：
\begin{lstlisting}
de_shell> read_file -format ddc MY_DESIGN.ddc
Reading ddc file '/designs/ex/MY_DESIGN.ddc'
Current design is 'MY_DESIGN'
\end{lstlisting}
  \item \cmd{elaborate} 命令
  \item \cmd{current\_design} 命令\\
        使用该命令将 \cmd{de\_shell} 内存中的任意设计设为当前设计。例如：
\begin{lstlisting}
de_shell> current_design MY_DESIGN
Current design is 'MY_DESIGN'. {MY_DESIGN}
\end{lstlisting}
\end{itemize}

要显示当前设计的名称，输入：
\begin{lstlisting}
de_shell> printvar current_design
current_design = "test"
\end{lstlisting}

为避免过度占用运行时间，避免编写包含大量 \cmd{current\_design} 命令的脚本（例如在循环中）。

% ============================================================
\section{链接设计}
\subsection*{Linking Designs}

要使设计完整，设计中的所有单元实例必须链接到所引用的库组件与设计。
要链接设计，使用 \cmd{link} 命令。
DC Explorer 使用由 \varname{link\_library} 变量、\varname{search\_path} 变量以及 \texttt{local\_link\_library} 属性设置的链接库来解析引用。

默认情况下，链接过程的大小写敏感性取决于引用的来源。
要设置链接过程的大小写敏感性，设置 \varname{link\_force\_case} 变量。

可执行的相关任务包括：
\begin{itemize}
  \item 更改设计引用
  \item 使用搜索路径定位设计
  \item 查询设计引用
  \item 解析引用的过程
\end{itemize}

\begin{seeAlsoBox}
\begin{itemize}
  \item 解析引用的过程
\end{itemize}
\end{seeAlsoBox}

\subsection{更改设计引用}
\subsubsection*{Changing Design References}

要更改单元或引用所链接到的组件或设计，使用 \cmd{change\_link} 命令。
当指定单元实例为对象时，更改该单元的链接；当指定引用为对象时，更改具有该引用的所有单元的链接。
仅当指定的 reference component 或 design 与原 reference 的端口数量相同，且对应端口的位宽和方向均相同时，DC Explorer 才会更改 link。

运行 \cmd{change\_link} 会将所有链接信息从旧设计复制到新设计。
若旧设计是 synthetic 模块，旧 synthetic 模块的所有属性会移到新链接。
运行 \cmd{change\_link} 后，必须对设计运行 \cmd{link} 命令。

可使用 \opt{-all\_instances} 选项指定位于层次较低层、且其父设计不唯一的实例作为对象。
同一父设计中的所有相似单元会自动链接到新引用设计；无需更改当前设计即可更改此类单元的链接。

\paragraph{示例 1}
下列命令将当前设计中 U1 与 U2 单元的链接更改为 MY\_ADDER 设计：
\begin{lstlisting}
de_shell> copy_design ADDER MY_ADDER
de_shell> change_link {U1 U2} MY_ADDER
\end{lstlisting}

\paragraph{示例 2}
下列命令更改位于层次较低层的 U1 单元的链接：
\begin{lstlisting}
de_shell> change_link top/sub_inst/U1 lsi_10k/AN3
\end{lstlisting}

\paragraph{示例 3}
下列命令在父设计 \texttt{bot} 被实例化两次（\texttt{mid1/bot1} 与 \texttt{mid1/bot2}）时，更改 \texttt{inv1} 单元所有实例的链接：
\begin{lstlisting}
de_shell> change_link -all_instances mid1/bot1/inv1 lsi_10k/AN3
Information: Changed link for all instances of cell 'inv1'
in subdesign 'bot'. (UID-193)

de_shell> get_cells -hierarchical -filter "ref_name == AN3"
{mid1/bot1/inv1 mid1/bot2/inv1}
\end{lstlisting}

\paragraph{示例 4}
可指示工具解析与原单元引用等效但引脚名不同的单元引用。
要在原引用与新单元引用之间映射引脚名，对 \cmd{change\_link} 使用
\opt{-pin\_map} \texttt{\{\{old\_pin1 new\_pin1\} \{old\_pin2 new\_pin2\}...\}} 选项。

例如，下列命令将当前设计中的 U1 单元链接到 AN21 单元，并把旧引脚名（A1、A2 与 Z）映射到新引脚名（A、B 与 Y）：
\begin{lstlisting}
de_shell> change_link [get_cells U1] AN21 \
   -pin_map {{A1 A} {A2 B} {Z Y}}
\end{lstlisting}

\subsection{使用搜索路径定位设计}
\subsubsection*{Locating Designs by Using a Search Path}

要查找设计文件位置，用 \cmd{which} 命令指定文件名。
DC Explorer 使用 \varname{search\_path} 变量定义的搜索路径查找设计文件，从该变量中最左侧目录开始，并报告找到的第一个设计文件。
默认情况下，搜索路径包括当前工作目录与 \texttt{\$SYNOPSYS/libraries/syn}，其中 \texttt{\$SYNOPSYS} 是安装目录路径。例如：
\begin{lstlisting}
de_shell> which my_design.ddc
{/usr/designers/example/my_design.ddc}
\end{lstlisting}

要在默认搜索路径之外指定其他目录，输入类似下列命令：
\begin{lstlisting}
de_shell> lappend search_path project
\end{lstlisting}

\subsection{查询设计引用}
\subsubsection*{Querying Design References}

可用下列方法查询设计引用：
\begin{itemize}
  \item 要报告当前实例或当前设计中所有引用的信息，使用 \cmd{report\_reference} 命令。
  \item 要返回您指定的实例集合，使用 \cmd{get\_references} 命令。\\
        例如，下列命令返回当前设计中具有 AN2 引用的实例集合：
\begin{lstlisting}
de_shell> get_references AN2
{U2 U3 U4}
\end{lstlisting}
  \item 要查看引用名称，使用 \cmd{report\_cell} 命令。例如：
\begin{lstlisting}
de_shell> report_cell [get_references AN*]
Cell                      Reference       Library         Area
Attributes
----------------------------------------------------------------
U2                        AN2             lsi_10k         2.000000
U3                        AN2             lsi_10k         2.000000
U4                        AN2             lsi_10k         2.000000
U8                        AN3             lsi_10k         2.000000
\end{lstlisting}
\end{itemize}

\subsection{解析引用的过程}
\subsubsection*{The Process of Resolving References}

DC Explorer 按下列步骤解析引用：
\begin{enumerate}
  \item 确定当前设计及其层次中引用了哪些库组件与子设计。
  \item 在链接库中搜索以定位引用。
  \begin{enumerate}
    \item DC Explorer 首先搜索当前设计的 \texttt{local\_link\_library} 属性中定义的库与设计文件。\\
          在层次化设计中，DC Explorer 仅考虑顶层设计的本地链接库，忽略与子设计关联的本地链接库。
    \item 若 \varname{link\_library} 变量设为星号（\texttt{*}），DC Explorer 在内存中搜索该引用。
    \item 然后 DC Explorer 搜索 \varname{link\_library} 变量中定义的库与设计文件。
  \end{enumerate}
  \item 若在链接库中未找到引用，则在 \varname{search\_path} 变量指定的目录中搜索。
  \item 将找到的引用链接（连接）到设计。
\end{enumerate}

DC Explorer 使用它定位到的第一个引用。
若定位到同名的其他引用，会生成警告消息标识被忽略的引用。
若未找到引用，则发出警告。

图~\ref{fig:dce-resolve-ref} 中的箭头显示链接过程在实例、引用与链接库之间创建的连接。
本例中，DC Explorer 在 LIBRARY\_2 逻辑库中找到 NAND2 组件，并在设计文件中找到 MULTIPLIER 子设计。

\figplaceholder{Figure 21: Resolving References}{解析引用}{fig:dce-resolve-ref}

\begin{seeAlsoBox}
\begin{itemize}
  \item 使用搜索路径定位设计
\end{itemize}
\end{seeAlsoBox}

% ============================================================
\section{编辑设计}
\subsection*{Editing Designs}

使用「设计编辑任务与命令」中描述的命令编辑设计。
对于 unique design，这些 netlist editing 命令接受 instance object，即层次中较低层的 cell。
可在层次化设计的任意层级操作，而无需使用 \cmd{current\_design} 命令。
例如，可输入下列命令在设计 \texttt{mid1} 中创建名为 \texttt{my\_cell} 的单元：
\begin{lstlisting}
de_shell> create_cell mid1/my_cell my_lib/AND2
\end{lstlisting}

连接或断开网络时，使用 \cmd{all\_connected} 命令查看连接到网络、端口或引脚的对象。
例如，下列命令序列用高功率反相器替换 U8 单元的引用：
\begin{lstlisting}
de_shell> get_pins U8/*
{"U8/A", "U8/Z"}
de_shell> all_connected U8/A
{"n66"}
de_shell> all_connected U8/Z
{"OUTBUS[10]"}
de_shell> remove_cell U8
Removing cell 'U8' in design 'top'.
1
de_shell> create_cell U8 IVP
Creating cell 'U8' in design 'top'.
1
de_shell> connect_net n66 [get_pins U8/A]
Connecting net 'n66' to pin 'U8/A'.
1
de_shell> connect_net OUTBUS[10] [get_pins U8/Z]
Connecting net 'OUTBUS[10]' to pin 'U8/Z'.
\end{lstlisting}

也可使用 \cmd{change\_link} 命令达到相同结果。例如：
\begin{lstlisting}
de_shell> change_link U8 IVP
\end{lstlisting}

其他网表编辑命令包括 \cmd{size\_cell}、\cmd{insert\_buffer} 与 \cmd{remove\_buffer}。

要查询或指定特定设计与对象，见：
\begin{itemize}
  \item 设计编辑任务与命令
  \item 查询设计与对象
  \item 指定设计对象
\end{itemize}

\begin{seeAlsoBox}
\begin{itemize}
  \item 设计编辑任务与命令
\end{itemize}
\end{seeAlsoBox}

\subsection{设计编辑任务与命令}
\subsubsection*{Design Editing Tasks and Commands}

可使用下表所列 \cmd{de\_shell} 命令对内存中的设计做增量编辑或更改网表。

\begin{table}[htbp]
\centering
\caption{设计编辑任务与命令}
\label{tab:dce-edit-cmds}
\begin{tabular}{@{}lll@{}}
\toprule
对象 & 任务 & 命令 \\
\midrule
Cells & 创建单元 & \cmd{create\_cell} \\
 & 删除单元 & \cmd{remove\_cell} \\
Nets & 创建网络 & \cmd{create\_net} \\
 & 连接网络 & \cmd{connect\_net} \\
 & 断开网络 & \cmd{disconnect\_net} \\
 & 删除网络 & \cmd{remove\_net} \\
Ports & 创建端口 & \cmd{create\_port} \\
 & 删除端口 & \cmd{remove\_port} \\
 &  & \cmd{remove\_unconnected\_ports} \\
Pins & 连接引脚 & \cmd{connect\_pin} \\
Buses & 创建总线 & \cmd{create\_bus} \\
 & 删除总线 & \cmd{remove\_bus} \\
\bottomrule
\end{tabular}
\end{table}

\begin{seeAlsoBox}
\begin{itemize}
  \item 编辑设计
\end{itemize}
\end{seeAlsoBox}

\subsection{查询设计与对象}
\subsubsection*{Querying the Design and Objects}

要查明内存中有哪些设计，使用 \cmd{list\_designs} 命令。
若还想查明对应的设计文件名，使用 \opt{-show\_file} 选项。例如：
\begin{lstlisting}
de_shell> list_designs
A (*)    B    C

de_shell> list_designs -show_file

/user1/designs/design_A/A.ddc
A (*)

/home/designer/dc/B.ddc
B       C
\end{lstlisting}

星号（\texttt{*}）标记当前设计。B 与 C 设计都包含在 \texttt{B.ddc} 文件中。

可使用表~\ref{tab:dce-query-objs} 所列命令列出当前设计中的各类对象。

\begin{table}[htbp]
\centering
\caption{查询设计对象的命令}
\label{tab:dce-query-objs}
\small
\begin{tabular}{@{}llp{6.5cm}@{}}
\toprule
对象 & 命令 & 任务 \\
\midrule
Instance & \cmd{list\_instances} & 列出实例及其引用。 \\
 & \cmd{report\_cell} & 显示实例信息。 \\
Reference & \cmd{report\_reference} & 显示引用信息。 \\
Port & \cmd{report\_port} & 显示端口信息。 \\
 & \cmd{report\_bus} & 显示总线端口信息。 \\
 & \cmd{all\_inputs} & 返回所有输入端口。 \\
 & \cmd{all\_outputs} & 返回所有输出端口。 \\
Net & \cmd{report\_net} & 显示网络信息。 \\
 & \cmd{report\_bus} & 显示总线网络信息。 \\
Clock & \cmd{report\_clock} & 显示时钟信息。 \\
 & \cmd{all\_clocks} & 返回所有时钟。 \\
Register & \cmd{all\_registers} & 返回所有寄存器。 \\
Collection & \cmd{get\_cells}、\cmd{get\_designs} 等 & 返回单元、设计、库与库单元、库单元引脚、网络、引脚或端口的集合。 \\
\bottomrule
\end{tabular}
\end{table}

\subsection{指定设计对象}
\subsubsection*{Specifying Design Objects}

可通过下列方式指定设计对象：
\begin{itemize}
  \item 使用相对路径
  \item 使用绝对路径
\end{itemize}

\subsubsection{使用相对路径}
\subsubsection*{Using a Relative Path}

使用相对路径指定设计对象时，对象必须位于当前设计中，且路径相对于当前实例。
当前实例是当前设计内的参照系。
默认情况下，当前实例是当前设计的顶层。
要更改当前实例，使用 \cmd{current\_instance} 命令。

例如，要在 Count\_16 设计中用 \texttt{dont\_touch} 属性标记 U1/U15 层次单元，使用下列命令序列之一：
\begin{itemize}
  \item 参照系保持在 Count\_16 设计顶层
\begin{lstlisting}
de_shell> current_design Count_16
Current design is 'Count_16'.
{Count_16}
de_shell> set_dont_touch U1/U15
\end{lstlisting}
  \item 参照系更改为 U1 实例
\begin{lstlisting}
de_shell> current_design Count_16
Current design is 'Count_16'.
{Count_16}
de_shell> current_instance U1
Current instance is '/Count_16/U1'.
/Count_16/U1
de_shell> set_dont_touch U15
\end{lstlisting}
\end{itemize}

此后 DC Explorer 将所有未来对象规范解释为相对于 U1 实例。
\cmd{current\_instance} 命令指向当前实例。

要显示当前实例，输入：
\begin{lstlisting}
de_shell> printvar current_instance
current_instance = "Count_16/U1"
\end{lstlisting}

要将当前实例复位到当前设计顶层，使用不带参数的 \cmd{current\_instance} 命令。例如：
\begin{lstlisting}
de_shell> current_instance
Current instance is the top-level of the design 'Count_16'
\end{lstlisting}

\subsubsection{使用绝对路径}
\subsubsection*{Using an Absolute Path}

使用绝对路径指定对象时，对象可以位于 \cmd{de\_shell} 内存中的任意设计。
指定对象的语法为：
\begin{lstlisting}
[file:]design/object
\end{lstlisting}
其中
\begin{itemize}
  \item \texttt{file} 是内存中设计文件的路径名，后跟冒号（\texttt{:}）\\
        当内存中多个设计同名时使用此参数。
  \item \texttt{design} 是 \cmd{de\_shell} 内存中的设计名
  \item \texttt{object} 是设计对象名，包括其层次路径
\end{itemize}

若不同类型的多个对象同名且未指定对象类型，DC Explorer 使用命令允许的类型查找对象。
要指定对象类型，使用 \cmd{get\_*} 命令。

例如，要在 Count\_16 设计中的 U1/U5 层次单元上放置 \texttt{dont\_touch} 属性，输入：
\begin{lstlisting}
de_shell> set_dont_touch /usr/designs/Count_16.ddc:Count_16/U1/U5
\end{lstlisting}

\begin{seeAlsoBox}
\begin{itemize}
  \item 查询设计与对象
\end{itemize}
\end{seeAlsoBox}

% ============================================================
\section{更改设计层次}
\subsection*{Changing the Design Hierarchy}

理想情况下，DC Explorer 中的设计层次应准确反映 HDL 描述中的设计划分。
不过，可在 DC Explorer 中更改层次，以便在不修改原始 HDL 描述的情况下试验替代层次划分。

在做任何设计层次更改之前，使用 \cmd{report\_hierarchy} 命令显示当前设计层次。
可用下列方法在不修改 HDL 描述的情况下更改层次。
\begin{itemize}
  \item 添加层次级别
  \item 移除层次级别
  \item 合并来自不同子设计的单元
\end{itemize}

做任何设计层次更改之后，使用 \cmd{report\_hierarchy} 命令验证层次更改。

\subsection{添加层次级别}
\subsubsection*{Adding Levels of Hierarchy}

要通过将单元或相关组件分组到子设计中来创建一层层次，使用 \cmd{group} 命令。
被分组的单元由引用新子设计的新实例替换。
新子设计的端口以其在设计中所连接的网络命名。
新子设计各端口的方向由对应网络的引脚确定。
对 cell 执行 group 操作时，原 cell 的属性和约束可能无法全部保留。

下列示例将单元分组到子设计中：

要将两个单元分组到名为 \texttt{my\_design}、实例名为 U 的新设计，输入：
\begin{lstlisting}
de_shell> group {u1 u2} -design_name my_design -cell_name U
\end{lstlisting}

要将所有具有 \texttt{alu} 前缀的单元分组到名为 \texttt{uP}、实例名为 UCELL 的新设计，输入：
\begin{lstlisting}
de_shell> group "alu*" -design_name uP -cell_name UCELL
\end{lstlisting}

要将具有唯一父设计 \texttt{mid1} 的 bot1、bot2 与 j 单元分组到名为 \texttt{my\_design}、实例名为 U1 的新子设计，输入：
\begin{lstlisting}
de_shell> group {mid1/bot1 mid1/bot2 mid1/j} \
   -cell_name U1 -design_name my_design
\end{lstlisting}

也可使用下列两条命令：
\begin{lstlisting}
de_shell> current_design mid1
de_shell> group {bot1 bot2 j} -cell_name U1 -design_name my_design
\end{lstlisting}

下列示例将相关组件分组到子设计中：

要将 process \texttt{ftj} 中 HDL 函数 \texttt{abc} 的所有单元分组到 \texttt{new\_block} 设计，输入：
\begin{lstlisting}
de_shell> group -hdl_block ftj/abc -design_name new_block
\end{lstlisting}

要将 process \texttt{ftj} 下所有总线化门分组到各自独立的层次级别，输入：
\begin{lstlisting}
de_shell> group -hdl_block ftj -hdl_bussed
\end{lstlisting}

\subsection{移除层次级别}
\subsubsection*{Removing Levels of Hierarchy}

DC Explorer 不会跨层次边界优化逻辑；因此，在某些设计中可能需要移除层次级别以改善时序。
移除一层层次称为 ungrouping，即将某一层次级别的子设计合并到父设计中。
标记了 \texttt{dont\_touch} 属性的设计、子设计与单元不能被 ungroup。

要 ungroup 层次，可使用下列命令：
\begin{itemize}
  \item 优化前的 \cmd{ungroup} 命令\\
        见「优化前 Ungroup 层次」。
  \item 优化期间显式使用的 \cmd{set\_ungroup} 命令\\
        见「优化期间显式 Ungroup 层次」。
  \item 优化期间隐式使用的 \cmd{compile\_exploration} 命令\\
        见「优化期间自动 Ungroup 层次」。
\end{itemize}

\begin{seeAlsoBox}
\begin{itemize}
  \item Ungroup 期间的时序约束保留
\end{itemize}
\end{seeAlsoBox}

\subsubsection{Ungroup 期间的时序约束保留}
\subsubsection*{Timing Constraint Preservation During Ungrouping}

单元被 ungroup 时，层次引脚会被移除。
DC Explorer 根据您是在优化前还是优化后 ungroup 层次，以不同方式处理放置在层次引脚上的时序约束。

表~\ref{tab:dce-ungroup-tc} 汇总不同 compile 流程中 ungrouping 对时序约束的影响。

\begin{table}[htbp]
\centering
\caption{保留层次引脚时序约束}
\label{tab:dce-ungroup-tc}
\small
\begin{tabular}{@{}p{6.2cm}p{7.5cm}@{}}
\toprule
Compile 流程 & 对层次引脚时序约束的影响 \\
\midrule
优化前使用 \cmd{ungroup} 命令 ungroup 层次 & 放置在层次引脚上的时序约束被保留。 \\
优化期间使用 \cmd{set\_ungroup} 后跟 \cmd{compile\_exploration} ungroup 层次 & 默认不保留层次引脚上的时序约束。要保留时序约束，将 \varname{auto\_ungroup\_preserve\_constraints} 变量设为 \texttt{true}。 \\
优化期间使用 \cmd{compile\_exploration} 自动 ungroup 层次 & 要 ungroup 层次并保留时序约束，将 \varname{auto\_ungroup\_preserve\_constraints} 变量设为 \texttt{true}。 \\
\bottomrule
\end{tabular}
\end{table}

保留时序约束时，DC Explorer 将时序约束重新赋给 ungroup 后仍留在同一网络上的相邻适当引脚。
约束会沿同一网络向前或向后移动。
注意：仅当驱动给定层次引脚的引脚不驱动其他引脚时，约束才能向后移动；否则约束必须向前移动。

当 DC Explorer 将约束移到叶单元时，通过用 \texttt{size\_only} 属性标记该单元，在 compile 期间保留约束。
当向前与向后方向都可行时，DC Explorer 选择使需标记 \texttt{size\_only} 属性的叶单元数量最少的方向。

对未映射设计做 ungroup 时，DC Explorer 将层次引脚上的约束移到叶单元并用 \texttt{size\_only} 属性标记该单元。
仅当未映射单元与目标库单元之间存在一对一匹配时，约束才会在 compile 过程中被保留。

仅保留用下列命令设置的时序约束：
\begin{lstlisting}
set_false_path
set_multicycle_path
set_min_delay
set_max_delay
set_input_delay
set_output_delay
set_disable_timing
set_case_analysis
create_clock
create_generated_clock
set_propagated_clock
set_clock_latency
\end{lstlisting}

\begin{noteBox}
\cmd{set\_rtl\_load} 约束不被保留。
此外，仅保留当前设计的时序约束；对当前设计层次做 ungroup 可能导致其他设计中的时序约束丢失。
\end{noteBox}

\subsubsection{优化前 Ungroup 层次}
\subsubsection*{Ungrouping Hierarchies Before Optimization}

使用 \cmd{ungroup} 命令在优化前 ungroup 一个或多个设计。
可在层次化设计的任意层级使用该命令，而无需使用 \cmd{current\_design} 命令。

\begin{noteBox}
若 ungroup 单元后使用 \cmd{change\_names} 命令修改层次分隔符（\texttt{/}），可能会丢失属性与约束信息。
\end{noteBox}

下列示例说明如何使用 \cmd{ungroup} 命令：

要 ungroup 单元列表，输入：
\begin{lstlisting}
de_shell> ungroup {high_decoder_cell low_decoder_cell}
\end{lstlisting}

要 ungroup U1 单元并为新单元指定 \texttt{U1:} 前缀，输入：
\begin{lstlisting}
de_shell> ungroup U1 -prefix "U1:"
\end{lstlisting}

要完全折叠当前设计的层次，输入：
\begin{lstlisting}
de_shell> ungroup -all -flatten
\end{lstlisting}

要递归 ungroup 属于 CELL\_X 单元、且位于当前设计下三层层次的单元，输入：
\begin{lstlisting}
de_shell> ungroup -start_level 3 CELL_X
\end{lstlisting}

要递归 ungroup 位于当前设计下三层层次、且属于 U1 与 U2 单元的单元（U1 与 U2 是当前设计的子单元），输入：
\begin{lstlisting}
de_shell> ungroup -start_level 3 {U1 U2}
\end{lstlisting}

要递归 ungroup 位于当前设计下三层层次的所有单元，输入：
\begin{lstlisting}
de_shell> ungroup -start_level 3 -all
\end{lstlisting}

实例对象是较低层次的单元。
下例说明当父设计唯一时，\cmd{ungroup} 如何接受实例对象。
本例中，MID1/BOT1 是设计 BOT 的唯一实例化。
该命令 ungroup MID1/BOT1/DES1 与 MID1/BOT1/DES2 单元。
\begin{lstlisting}
de_shell> ungroup {MID1/BOT1/DES1 MID1/BOT1/DES2}
\end{lstlisting}

也可使用下列两条命令：
\begin{lstlisting}
de_shell> current_instance MID1/BOT1
de_shell> ungroup {DES1 DES2}
\end{lstlisting}

\begin{seeAlsoBox}
\begin{itemize}
  \item Ungroup 期间的时序约束保留
\end{itemize}
\end{seeAlsoBox}

\subsubsection{优化期间显式 Ungroup 层次}
\subsubsection*{Ungrouping Hierarchies Explicitly During Optimization}

可通过 \cmd{set\_ungroup} 后跟 \cmd{compile\_exploration} 控制优化期间要 ungroup 的设计。
指定要 ungroup 的单元或设计，\cmd{set\_ungroup} 会将 \texttt{ungroup} 属性赋给指定单元或被引用的设计。
若在设计上设置该属性，则引用该设计的所有单元都会被 ungroup。

例如，要在优化期间 ungroup U1 单元，输入：
\begin{lstlisting}
de_shell> set_ungroup U1
de_shell> compile_exploration
\end{lstlisting}

要查看对象是否设置了 \texttt{ungroup} 属性，使用 \cmd{get\_attribute} 命令：
\begin{lstlisting}
de_shell> get_attribute object ungroup
\end{lstlisting}

要移除 \texttt{ungroup} 属性，使用 \cmd{remove\_attribute} 命令，或通过 \cmd{set\_ungroup} 将 \texttt{ungroup} 属性设为 \texttt{false}。例如：
\begin{lstlisting}
de_shell> set_ungroup object false
\end{lstlisting}

\begin{seeAlsoBox}
\begin{itemize}
  \item Ungroup 期间的时序约束保留
\end{itemize}
\end{seeAlsoBox}

\subsubsection{优化期间自动 Ungroup 层次}
\subsubsection*{Ungrouping Hierarchies Automatically During Optimization}

要在优化期间自动 ungroup 层次，可使用 \cmd{compile\_exploration} 命令。
DC Explorer 提供两种自动 ungroup 策略：基于面积的自动 ungrouping 与基于延时的自动 ungrouping。

默认情况下，\cmd{compile\_exploration} 执行基于延时的自动 ungrouping，它沿关键路径 ungroup 层次，主要用于时序优化。
此外，该命令在初始映射之前执行基于面积的自动 ungrouping。
DC Explorer 估计未映射层次的面积，并移除较小的子设计，以改善面积与时序 QoR。

要报告在 \cmd{compile\_exploration} 过程中被 ungroup 的层次，使用 \cmd{report\_auto\_ungroup} 命令。
该报告给出每个被 ungroup 层次的实例名、单元名与实例数。

更多信息亦见 \textit{Design Compiler User Guide}。

\begin{seeAlsoBox}
\begin{itemize}
  \item Ungroup 期间的时序约束保留
\end{itemize}
\end{seeAlsoBox}

\subsection{合并来自不同子设计的单元}
\subsubsection*{Merging Cells From Different Subdesigns}

要将不同子设计中的单元合并到新子设计中：
\begin{enumerate}
  \item 将单元 group 到新设计中。
  \item Ungroup 该新设计。
\end{enumerate}

例如，下列命令序列创建包含最初位于 \texttt{u\_add} 与 \texttt{u\_mult} 子设计中之单元的新 \texttt{alu} 设计。
\begin{lstlisting}
de_shell> group {u_add u_mult} -design alu
de_shell> current_design alu
de_shell> ungroup -all
de_shell> current_design top_design
\end{lstlisting}

% ============================================================
\section{使用设计对象属性}
\subsection*{Using Design Object Attributes}

attribute 描述 design database 中对象的逻辑、电气、物理及其他特性。
attribute 附着于 design object，并在写出 design database 时随之保存；但修改 attribute 设置后，database 不会自动写回。

DC Explorer 在下列类型的对象上使用属性：
\begin{itemize}
  \item 整个设计
  \item 设计对象，如时钟、网络、引脚与端口
  \item 设计内的设计引用与单元实例
  \item 逻辑库、库单元与单元引脚
\end{itemize}

属性有名称、类型与值。属性类型可以是：字符串、数值或逻辑（Boolean）。

某些属性为只读，由 DC Explorer 设置这些属性值。
不能更改只读属性的值，但可更改读/写属性的值。

要了解如何查询、更改、在设计对象上设置属性，见：
\begin{itemize}
  \item 查询属性值
  \item 在对象上设置属性
  \item 更改对象属性
  \item 获取属性描述的命令
  \item 传播在单元实例上设置的用户属性
  \item 在时序单元上保留用户属性
  \item 对象搜索顺序
\end{itemize}

\subsection{查询属性值}
\subsubsection*{Querying Attribute Values}

要查询设计对象属性的值，使用 \cmd{report\_attribute} 命令、\cmd{get\_attribute} 命令或 GUI。

\paragraph{所有属性}
\paragraph*{All Attributes}

使用 \cmd{report\_attribute} 命令查看所有属性。例如：
\begin{lstlisting}
de_shell> report_attribute object_list
\end{lstlisting}

\paragraph{特定属性}
\paragraph*{Specific Attributes}

使用 \cmd{get\_attribute} 命令查看特定属性。下例报告 OUT7 端口上的 \texttt{max\_fanout} 属性：
\begin{lstlisting}
de_shell> get_attribute OUT7 max_fanout
Performing get_attribute on port 'OUT7'.
{3.000000}
\end{lstlisting}

可查询引脚与端口的边界框（bounding box）信息。例如，下列命令报告寄存器 \texttt{reg\_1} 的 Q 引脚与当前设计的 \texttt{test\_mode} 端口上的 \texttt{bbox} 属性：
\begin{lstlisting}
de_shell> get_attribute [get_pins reg_1/Q] bbox
de_shell> get_attribute [get_ports test_mode] bbox
\end{lstlisting}

\begin{noteBox}
要查询引脚与端口的 \texttt{bbox} 属性，需通过将 \varname{de\_enable\_physical\_flow} 变量设为 \texttt{true} 来启用物理流程能力。
\end{noteBox}

若某属性适用于多种对象类型，DC Explorer 在数据库中搜索指定对象。

\paragraph{使用 GUI}
\paragraph*{Using the GUI}

可在 GUI 的 Properties 对话框中查看所选设计、设计对象或时序路径的属性及其他对象属性。
也可为某些属性设置、更改或移除属性值。

有关使用 GUI 的更多信息，见 \textit{Design Vision User Guide}。

\begin{seeAlsoBox}
\begin{itemize}
  \item 在对象上设置属性
  \item 更改对象属性
  \item 获取属性描述的命令
  \item 对象搜索顺序
\end{itemize}
\end{seeAlsoBox}

\subsection{在对象上设置属性}
\subsubsection*{Setting Attributes on Objects}

要在设计对象上设置属性，使用下列两种方法之一：
\begin{itemize}
  \item 使用属性专用命令\\
        使用属性专用命令在对象上设置与该命令关联的属性。下例用 \texttt{dont\_touch} 属性标记 U1 单元。
\begin{lstlisting}
de_shell> set_dont_touch U1
\end{lstlisting}
  \item 使用 \cmd{set\_attribute} 命令\\
        使用该命令设置任意属性的值，或定义新属性并设置其值。
        该命令强制执行预定义属性类型（如字符串或整数）。
        若以不正确类型的值设置属性，命令会发出错误消息。
        要查找属性的预定义类型，使用 \cmd{list\_attributes} 命令。\\
        在引用（子设计或库单元）上设置属性时，该属性适用于设计中具有该引用的所有单元。
        在实例（单元、网络或引脚）上设置属性时，该属性覆盖从实例引用继承的任何属性。\\
        下例在 \texttt{lsi\_10k/FJK3} 库单元上设置 \texttt{dont\_touch} 属性：
\begin{lstlisting}
de_shell> set_attribute lsi_10K/FJK3 dont_touch true
\end{lstlisting}
\end{itemize}

\begin{seeAlsoBox}
\begin{itemize}
  \item 查询属性值
  \item 更改对象属性
  \item 获取属性描述的命令
  \item 对象搜索顺序
\end{itemize}
\end{seeAlsoBox}

\subsection{更改对象属性}
\subsubsection*{Changing Object Attributes}

可更改、移除或保存对象上设置的属性：
\begin{itemize}
  \item 使用 \cmd{remove\_attribute} 命令移除对象上的特定属性\\
        不能用该命令移除继承的属性。例如，若 \texttt{dont\_touch} 属性赋给引用，应从引用而非从继承该属性的单元移除该属性。\\
        下例从 OUT7 端口移除 \texttt{max\_fanout} 属性：
\begin{lstlisting}
de_shell> remove_attribute OUT7 max_fanout
\end{lstlisting}
  \item 使用 \cmd{reset\_design} 命令从当前设计移除所有属性\\
\cmd{reset\_design} 移除所有 design information，包括 clock、input/output delay、path group、operating condition、timing range 与 wire load model。
        运行该命令通常等效于重新开始设计。
  \item 更改属性值\\
        要更改属性值，先移除属性，再重新创建以设置类型。
  \item 使用 \cmd{write\_script} 命令保存属性\\
        退出 \cmd{de\_shell} 时，DC Explorer 不会自动保存属性值。
        要重新创建属性，使用 \cmd{write\_script} 生成 \cmd{de\_shell} 脚本。
        默认情况下，\cmd{write\_script} 将属性设置命令报告到屏幕。
        使用重定向运算符（\texttt{>}）将输出重定向到文件。例如：
\begin{lstlisting}
de_shell> write_script > attr.scr
\end{lstlisting}
\end{itemize}

\begin{noteBox}
\cmd{write\_script} 命令仅支持应用（系统定义）属性，不支持用户定义属性。
\end{noteBox}

\begin{seeAlsoBox}
\begin{itemize}
  \item 查询属性值
  \item 更改对象属性
  \item 获取属性描述的命令
  \item 对象搜索顺序
\end{itemize}
\end{seeAlsoBox}

\subsection{获取属性描述的命令}
\subsubsection*{Commands to Get Attribute Descriptions}

多数属性仅适用于有限选择的对象类型；例如，\texttt{rise\_drive} 属性仅适用于输入与 inout 端口。
某些属性适用于多种对象类型；例如，\texttt{dont\_touch} 属性可适用于网络、单元、端口、引用或设计。
某些属性预定义并由 DC Explorer 识别；其他属性由用户定义。
可使用下表所列命令获取预定义属性的详细信息。

\begin{table}[htbp]
\centering
\caption{获取属性描述的命令}
\label{tab:dce-attr-desc}
\begin{tabular}{@{}ll@{}}
\toprule
对象类型 & 命令 \\
\midrule
All & \texttt{man attributes} \\
Designs & \texttt{man design\_attributes} \\
Cells & \texttt{man cell\_attributes} \\
Clocks & \texttt{man clock\_attributes} \\
Nets & \texttt{man net\_attributes} \\
Pins & \texttt{man pin\_attributes} \\
Ports & \texttt{man port\_attributes} \\
Libraries & \texttt{man library\_attributes} \\
Library cells & \texttt{man library\_cell\_attributes} \\
References & \texttt{man reference\_attributes} \\
\bottomrule
\end{tabular}
\end{table}

\begin{seeAlsoBox}
\begin{itemize}
  \item 查询属性值
  \item 在对象上设置属性
  \item 更改对象属性
  \item 对象搜索顺序
\end{itemize}
\end{seeAlsoBox}

\subsection{传播在单元实例上设置的用户属性}
\subsubsection*{Propagating User Attributes Set on Cell Instances}

要将较低层设计单元实例上设置的用户定义属性传播到较高层层次中的单元实例，使用 \cmd{propagate\_user\_attributes} 命令。
例如，在图~\ref{fig:dce-cell-u1} 中：
\begin{itemize}
  \item 当前设计设为 \texttt{mid}，用户定义属性定义为 XYZ，在 \texttt{mid} 设计的所有单元上值为 \texttt{false}。
  \item 用户定义属性设置在 \texttt{bot} 设计的 B1 与 B2 实例上，这些实例包含单元实例 U1。
\end{itemize}

\figplaceholder{Figure 22: Cell Instances U1 in the bot Design}{\texttt{bot} 设计中的单元实例 U1}{fig:dce-cell-u1}

当当前设计（\texttt{mid}）更改为 Top 时（如示例~2 所示），使用 \cmd{propagate\_user\_attributes} 命令将 \texttt{bot} 设计中所有单元实例上设置的所有用户定义属性传播到 Top 设计。

\paragraph{示例 2}
\paragraph*{Example 2: Cell Instances U1 in the bot Design Propagate to the Top Design}
\begin{lstlisting}
prompt> current_design mid
Current design is 'mid'.
{mid}

prompt> define_user_attribute -type boolean -classes cell XYZ
XYZ

prompt> set_user_attribute [get_cells -hierarchical *] XYZ false
Information: Attribute 'XYZ' is set on 4 objects. (UID-186)
B1/U1 B2/U1 B1 B2

prompt> current_design Top
Current design is 'top'.
{top}

prompt> get_cells -hierarchical -filter "XYZ==false"
{M1/B1 M1/B2 M2/B1 M2/B2}

prompt> propagate_user_attributes {XYZ}
1

prompt> get_cells -hierarchical -filter "XYZ==false"
{M1/B1 M1/B1/U1 M1/B2 M1/B2/U1 M2/B1 M2/B1/U1 M2/B2 M2/B2/U1}
\end{lstlisting}

\begin{seeAlsoBox}
\begin{itemize}
  \item 在对象上设置属性
\end{itemize}
\end{seeAlsoBox}

\subsection{在时序单元上保留用户属性}
\subsubsection*{Preserving User Attributes on Sequential Cells}

要指定 synthesis flow 中应在 sequential cell 上保留哪些 user-defined attribute，使用 \cmd{define\_preserve\_user\_attribute} 命令。

\begin{noteBox}
\begin{itemize}
  \item 属性列表仅在当前会话中保留。
  \item 若在同一会话中多次使用该命令，它会替换先前保留的属性列表。
\end{itemize}
\end{noteBox}

要报告保留的属性列表，使用 \cmd{report\_preserve\_user\_attribute} 命令。

下例说明如何指定要保留的属性并列出已指定属性：
\begin{lstlisting}
de_shell> define_preserve_user_attribute {my_attribute_1 my_attribute_2}
1

de_shell> report_preserve_user_attribute
my_attribute_1 my_attribute_2
\end{lstlisting}

该功能有下列限制：
\begin{itemize}
  \item 只能为时序单元保留属性，不能为引脚或网络保留。
  \item \cmd{optimize\_registers} 与 \cmd{write\_script} 命令不支持 \cmd{define\_preserve\_user\_attribute} 命令。
\end{itemize}

\subsection{对象搜索顺序}
\subsubsection*{Object Search Order}

可在多种对象类型上设置属性的命令使用特定搜索顺序，以找到属性所适用的对象。
例如，\cmd{set\_dont\_touch} 命令作用于单元、网络、引用与库单元。
若对 \cmd{set\_dont\_touch} 指定名为 X 的对象，且有两个对象（如设计与单元）都名为 X，DC Explorer 将属性应用于找到的第一个对象类型。
此时属性设置在设计上，而非单元上。

DC Explorer 一直搜索直到找到匹配对象；若未找到匹配对象则显示错误消息。

可使用 \cmd{get\_*} 命令指定对象类型来搜索指定类型的对象。
例如，假定当前设计包含同名为 \texttt{my\_object} 的单元与网络。
下列命令因默认搜索顺序而在单元上设置 \texttt{dont\_touch} 属性：
\begin{lstlisting}
de_shell> set_dont_touch my_object
\end{lstlisting}

要改为在网络上设置 \texttt{dont\_touch} 属性，输入：
\begin{lstlisting}
de_shell> set_dont_touch [get_nets my_object]
\end{lstlisting}

\begin{seeAlsoBox}
\begin{itemize}
  \item 查询属性值
  \item 在对象上设置属性
  \item 更改对象属性
  \item 获取属性描述的命令
\end{itemize}
\end{seeAlsoBox}

% ============================================================
\section{创建设计}
\subsection*{Creating Designs}

要创建新设计，使用 \cmd{create\_design} 命令。
使用该命令在 \cmd{de\_shell} 内存中创建空设计。
下例创建名为 \texttt{my\_design} 的新设计。内存文件名为 \texttt{my\_design.db}，路径为当前工作目录。
\begin{lstlisting}
de_shell> create_design my_design
de_shell> list_designs -show_file

/work_dir/mapped/test.ddc
test (*) test_DW01_inc_16_0 test_DW02_mult_16_16_1

/work_dir/my_design.db
my_design
\end{lstlisting}

用 \cmd{create\_design} 创建的设计不包含设计对象。
要向新设计添加设计对象，使用如 \cmd{create\_clock}、\cmd{create\_cell}、\cmd{create\_bus}、\cmd{create\_port} 与 \cmd{connect\_net} 等命令。

\begin{seeAlsoBox}
\begin{itemize}
  \item 编辑设计
\end{itemize}
\end{seeAlsoBox}

% ============================================================
\section{复制设计}
\subsection*{Copying Designs}

使用 \cmd{copy\_design} 命令将内存中的设计复制为新设计。
下例复制内存中的 \texttt{test} 设计，并将新设计命名为 \texttt{test\_new}：
\begin{lstlisting}
de_shell> copy_design test test_new
Information: Copying design /designs/test.ddc:to
designs/test.ddc:test_new

de_shell> list_designs -show_file
/designs/test.ddc
test (*) test_new
\end{lstlisting}

要手动创建可编辑的新子设计，并将实例链接到这些新子设计，使用 \cmd{copy\_design} 与 \cmd{change\_link} 命令。
例如，假定设计有两个相同单元 U1 与 U2，都链接到 COMP。输入下列命令序列以创建新的、各自独立链接的子设计：
\begin{lstlisting}
de_shell> copy_design COMP COMP1
Information: Copying design /designs/COMP.ddc:COMP to
   designs/COMP.ddc:COMP1

de_shell> change_link U1 COMP1
Performing change_link on cell 'U1'.

de_shell> copy_design COMP COMP2
Information: Copying design /designs/COMP.ddc:COMP to
   designs/COMP.ddc:COMP2

de_shell> change_link U2 COMP2
Performing change_link on cell 'U2'.
\end{lstlisting}

% ============================================================
\section{重命名设计}
\subsection*{Renaming Designs}

可为设计指定新名称，或将设计列表移到文件。
要重命名内存中的设计，使用 \cmd{rename\_design} 命令。
要保存重命名后的文件，使用 \cmd{write\_file} 命令。
可使用 \cmd{list\_designs} 命令显示运行 \cmd{rename\_design} 前后的设计。例如：
\begin{lstlisting}
de_shell> list_designs -show_file
/designs/test.ddc
test(*) test_new

de_shell> rename_design test_new test_new_1
Information: Renaming design /designs/test.ddc:test_new to
/designs/test.ddc:test_new_1

de_shell> list_designs -show_file
/designs/test.ddc
test (*) test_new_1
\end{lstlisting}

可使用 \opt{-prefix}、\opt{-postfix} 与 \opt{-update\_links} 选项重命名设计，并为整个设计层次更新单元链接。
例如，下列命令序列为 D 设计添加 \texttt{NEW\_} 前缀，并更新其实例单元的链接：
\begin{lstlisting}
de_shell> get_cells -hierarchical -filter "ref_name == D"
{b_in_a/c_in_b/d1_in_c b_in_a/c_in_b/d2_in_c}

de_shell> rename_design D -prefix NEW_ -update_links
Information: Renaming design /test_dir/D.ddc:D to
/test_dir/D.ddc:NEW_D. (UIMG-45)

de_shell> get_cells -hierarchical -filter "ref_name == D"
# no such cells!

de_shell> get_cells -hierarchical -filter "ref_name == NEW_D"
{b_in_a/c_in_b/d1_in_c b_in_a/c_in_b/d2_in_c}
\end{lstlisting}

本例中，第一步报告实例化 D 设计的单元。运行第二步后，实例链接到重命名后的引用设计 NEW\_D。

% ============================================================
\section{保存设计}
\subsection*{Saving Designs}

可随时以不同名称或格式保存（写入磁盘）设计及其层次中的子设计。
修改设计后应手动保存，因为 DC Explorer 不会自动保存设计。

DC Explorer 支持下表所列设计文件格式。

\begin{table}[htbp]
\centering
\caption{支持的输出格式}
\label{tab:dce-output-fmt}
\small
\begin{tabular}{@{}lp{10cm}@{}}
\toprule
格式 & 说明 \\
\midrule
\texttt{.ddc} & Synopsys 内部数据库格式 \\
Verilog & IEEE 标准 Verilog（见 HDL Compiler 文档） \\
svsim & SystemVerilog 网表 wrapper。注意：\cmd{write\_file} \opt{-format} \texttt{svsim} 仅写出网表 wrapper，而不写出被测门级设计（DUT）本身。要写出门级 DUT，必须使用现有的 \cmd{write\_file} \opt{-format} \texttt{verilog} 命令。 \\
VHDL & IEEE 标准 VHDL（见 HDL Compiler 文档） \\
Milkyway & 在 DC Explorer 中写出 Milkyway 数据库的格式 \\
\bottomrule
\end{tabular}
\end{table}

可用下列方法保存设计文件：
\begin{itemize}
  \item 使用 \cmd{write\_file} 命令将内存中的设计转换为指定格式并保存到磁盘\\
        默认情况下，\cmd{write\_file} 仅保存顶层设计。要保存整个设计，指定 \opt{-hierarchy} 选项。
        下例使用 \opt{-hierarchy} 与 \opt{-format} 选项将 top 设计层次中的所有设计写入 \texttt{.ddc} 文件：
\begin{lstlisting}
de_shell> write_file -hierarchy -format ddc top
Writing ddc file 'top.ddc' Writing ddc file 'A.ddc' Writing ddc file
'B.ddc'
\end{lstlisting}
        下例使用 \opt{-output} 选项将多个设计写入 \texttt{test.ddc} 文件：
\begin{lstlisting}
de_shell> write_file -format ddc -output test.ddc {ADDer MULT16}
Writing ddc file 'test.ddc'
\end{lstlisting}
  \item 使用 \cmd{write\_milkyway} 命令写入 Milkyway 数据库\\
        \cmd{write\_milkyway} 基于内存中的网表创建设计文件，并将当前设计的设计数据保存在该文件中。
\end{itemize}

\begin{seeAlsoBox}
\begin{itemize}
  \item 设计数据库与网表命名一致性
  \item 重命名设计
  \item 使用 Milkyway 数据库
\end{itemize}
\end{seeAlsoBox}

% ============================================================
\section{移除设计}
\subsection*{Removing Designs}

完成 compile 会话并保存优化后的设计后，可在打开另一设计之前从内存中删除该设计。
要从内存移除设计，使用 \cmd{remove\_design} 命令。
默认情况下，\cmd{remove\_design} 仅从 \cmd{de\_shell} 内存中移除指定的 design。

若定义了引用设计对象的变量，从内存移除设计时 DC Explorer 会移除这些引用。
这可防止后续命令对不存在的设计对象操作。

例如：
\begin{lstlisting}
de_shell> set_app_var PORTS [all_inputs]
{"A0", "A1", "A2", "A3"}
de_shell> query_objects $PORTS
PORTS = {"A0", "A1", "A2", "A3"}
de_shell> remove_design
Removing design 'top'
1
de_shell> query_objects $PORTS
Error: No such collection '_sel2' (SEL-001)
\end{lstlisting}

% ============================================================
\section{设计数据库与网表命名一致性}
\subsection*{Design Database and Netlist Naming Consistency}

在 \cmd{de\_shell} 中写出网表之前，确保所有网络与端口名称符合布局工具的命名约定，并使用一致的总线命名风格。

某些 ASIC 与 EDA 厂商提供会创建 \texttt{.synopsys\_dc.setup} 文件的程序，其中包含将名称转换为其约定的适当命令。
若需要更改任何网络或端口名称，使用 \cmd{define\_name\_rules} 与 \cmd{change\_names} 命令。

本节主题：
\begin{itemize}
  \item \texttt{.synopsys\_dc.setup} 文件中的命名规则
  \item 解决命名问题
  \item 避免 SystemVerilog 与 VHDL 结构中的 bit-blasted 端口
\end{itemize}

\subsection{\texorpdfstring{\texttt{.synopsys\_dc.setup} 文件中的命名规则}{Naming Rules in the .synopsys\_dc.setup File}}
\subsubsection*{Naming Rules in the .synopsys\_dc.setup File}

使用 \cmd{define\_name\_rules} 命令定义设计对象命名规则集。
\texttt{.synopsys\_dc.setup} 文件显示特定布局工具厂商提供的下列命名规则。这些命名规则：
\begin{itemize}
  \item 将对象名限制为字母数字字符
  \item 将 DesignWare 单元名更改为有效名称（将 ``*cell*'' 改为 ``U''，将 ``*-return'' 改为 ``RET''）
\end{itemize}
\begin{lstlisting}
define_name_rules simple_names -allowed "A-Za-z0-9_" \
-last_restricted "_" \
-first_restricted "_" \
-map { {{"\*cell\*","U"}, {"*-return","RET"}} }
\end{lstlisting}

例如，可指定 \opt{-map} 选项以移除单元名的尾部下划线：
\begin{lstlisting}
de_shell> define_name_rules naming_convention \
   -map {{{_$, ""}} } -type cell
\end{lstlisting}

要列出可用命名规则，使用 \cmd{report\_name\_rules} 命令。

\begin{noteBox}
您的厂商可能使用不同命名约定。请与厂商确认需遵循的命名约定。
\end{noteBox}

\subsection{解决命名问题}
\subsubsection*{Resolving Naming Problems}

您可能在对象、输入输出文件与工具集中遇到命名约定冲突。
在设计数据库文件中，可有许多设计对象（如端口、网络、单元、逻辑模块与逻辑模块引脚），各自有命名约定。
此外，可能使用多种具有不同语法定义的输入输出文件格式。
使用多家厂商的工具集可能引入更多命名问题。

正确命名有助于消除不匹配错误，并减少设计中对名称转义（name escaping）的需要。
要解决命名问题，在写出任何文件之前使用 \cmd{change\_names} 命令解决命名问题并在设计数据库文件中进行名称更改。初始流程为：
\begin{enumerate}
  \item 读入设计 RTL 并施加约束。\\
        此处无需更改您的方法。
  \item 综合设计以生成门级描述。\\
        按通常方式再次 compile 或优化设计，使用标准脚本集。
  \item 应用名称更改并解决命名问题。在写出设计之前使用 \cmd{change\_names} 及其 Verilog 或 VHDL 开关。

\begin{noteBox}
每当要写出 Verilog 或 VHDL design 时，都应使用
\cmd{change\_names} \opt{-rules} \opt{-[verilog|vhdl]} \opt{-hierarchy} 命令，
因为 design database 中的命名不符合 Verilog 或 VHDL 语法。例如：
\begin{lstlisting}
de_shell> change_names -rules verilog -hierarchy
\end{lstlisting}
\end{noteBox}
  \item 使用 \cmd{write} \opt{-format} \texttt{verilog} 命令将文件写入磁盘。\\
        留意报告的名称更改，这表示需要重复步骤 3 并细化命名规则。
  \item 若所有适当的名称更改都已完成，输出文件与设计数据库文件匹配。输入下列命令并比较输出。
\begin{lstlisting}
de_shell> write -format verilog -hierarchy -output "consistent.v"
de_shell> write -format ddc -hierarchy -output "consistent.ddc"
\end{lstlisting}
  \item 为第三方工具写出文件。
\end{enumerate}

要在实际不进行更改的情况下显示 \cmd{change\_names} 命令的效果，使用 \cmd{report\_names} 命令。

\subsection{避免 SystemVerilog 与 VHDL 结构中的 bit-blasted 端口}
\subsubsection*{Avoiding Bit-Blasted Ports in SystemVerilog and VHDL Structures}

可使用 \cmd{change\_names} 命令避免 SystemVerilog struct 与 VHDL record（packed 或 unpacked）中的 bit-blasted 端口。
要启用此能力，对 \cmd{define\_name\_rules} 指定 \opt{-preserve\_struct\_ports} 选项。
可按如下方式将该选项添加到现有 Verilog 命名规则：
\begin{lstlisting}
de_shell> define_name_rules verilog -preserve_struct_ports
de_shell> change_names -hierarchy -rules verilog
\end{lstlisting}

注意：使用 \cmd{define\_name\_rules} 的下列任一选项会覆盖 \opt{-preserve\_struct\_ports} 选项：
\begin{itemize}
  \item \opt{-remove\_port\_bus}
  \item \opt{-dont\_change\_bus\_members}
\end{itemize}

如下列示例所示，Verilog 网表 2 使用前述两条命令指定的命名规则保留了 struct 端口：
\begin{itemize}
  \item RTL
\begin{lstlisting}
typedef struct {logic blue; logic red;} color;
module top (
   input  color i_color,
   output color o_color
);
assign o_color = i_color;
endmodule
\end{lstlisting}
  \item 具有 bit-blasted struct 端口的网表 1
\begin{lstlisting}
module top ( i_color_blue_, i_color_red_, o_color_blue_,
o_color_red_ );
   input  i_color_blue_, i_color_red_;
   output o_color_blue_, o_color_red_;
   wire   o_color_blue_, o_color_red_;
   assign o_color_blue_ = i_color_blue_;
   assign o_color_red_  = i_color_red_;
endmodule
\end{lstlisting}
  \item 保留 struct 端口的网表 2
\begin{lstlisting}
module top ( i_color, o_color );
   input  [1:0] i_color;
   output [1:0] o_color;
assign o_color[1] = i_color[1];
assign o_color[0] = i_color[0];
endmodule
\end{lstlisting}
\end{itemize}
